% ==============================================================================
% CAPÍTULO 4 - REVISÃO BIBLIOGRÁFICA
% ==============================================================================

\chapter{Revisão Bibliográfica}
\label{ch:revisao-bibliografica}

Este capítulo apresenta os fundamentos teóricos e tecnológicos que embasam o trabalho desenvolvido
durante o estágio. São abordados os principais conceitos relacionados ao desenvolvimento web com Django,
à geração de documentos com LaTeX e à modelagem de dados relacional, pilares sobre os quais a solução
de geração automatizada de propostas comerciais foi construída.

%-------------------------------------------------
\section{Django e o Padrão MTV}
%-------------------------------------------------
\label{sec:django}

Django é um \textit{framework} web de alto nível escrito em Python, projetado para facilitar o
desenvolvimento rápido de aplicações web seguras e de fácil manutenção \cite{django2024}. Sua arquitetura
segue o padrão MTV (\textit{Model-Template-View}), uma variação do amplamente conhecido MVC
(\textit{Model-View-Controller}), no qual o \textit{Model} define a estrutura dos dados e sua relação
com o banco de dados; o \textit{Template} é responsável pela apresentação, seja HTML, seja um documento
LaTeX, como utilizado neste projeto; e a \textit{View} concentra a lógica de negócio, recebendo
requisições HTTP, consultando os modelos e retornando uma resposta ao cliente \cite{holovaty2009}.

Um dos componentes centrais do Django utilizado neste projeto é o ORM (\textit{Object-Relational
Mapper}), que permite interagir com o banco de dados por meio de classes Python, abstraindo a escrita
de SQL. Cada modelo Django é uma subclasse de \texttt{django.db.models.Model}, e seus atributos
correspondem às colunas de uma tabela relacional. As migrações, geradas automaticamente pelo comando
\texttt{makemigrations}, traduzem alterações nos modelos em operações de banco de dados reproduzíveis e
controladas por versão \cite{django2024}.

O sistema de \textit{templates} do Django suporta uma linguagem própria de marcação que permite
interpolação de variáveis, estruturas de controle (\textit{if}, \textit{for}) e herança de templates.
Por padrão, os delimitadores dessa linguagem são chaves duplas para variáveis e blocos
percentual-chave para \textit{tags}, o que pode gerar conflitos com a sintaxe do LaTeX, cujo
caractere de percentual tem função de comentário e cujas chaves são usadas extensivamente como
delimitadores de comandos. Esse conflito demanda atenção durante o desenvolvimento de templates LaTeX
processados pelo Django.

%-------------------------------------------------
\section{LaTeX e Geração de Documentos}
%-------------------------------------------------
\label{sec:latex}

LaTeX é um sistema de preparação de documentos baseado no programa de composição tipográfica TeX,
desenvolvido por Donald Knuth \cite{knuth1984}. Diferentemente de editores \textit{WYSIWYG}
(\textit{What You See Is What You Get}), no LaTeX o autor descreve a estrutura e o conteúdo do documento
por meio de comandos textuais, e o sistema se encarrega da diagramação final, resultando em documentos
com qualidade tipográfica superior \cite{lamport1994}.

Para a compilação de documentos LaTeX em arquivos PDF, este projeto utiliza o compilador
\textbf{LuaLaTeX}, variante moderna do LaTeX com suporte nativo a Unicode e a fontes OpenType, em
conjunto com o utilitário \texttt{latexmk}, que automatiza o processo de compilação, executando o
compilador quantas vezes forem necessárias para resolver referências cruzadas, sumários e outros
elementos que dependem de múltiplas passagens \cite{latex2023}.

A abordagem adotada neste projeto, renderizar um template \texttt{.tex} com o sistema de templates do
Django e, em seguida, compilá-lo com \texttt{latexmk}, é uma estratégia consolidada para geração de
PDFs com alto controle tipográfico em aplicações web, especialmente quando o documento exige layout
complexo, como tabelas posicionadas com precisão, imagens integradas e fontes customizadas.

%-------------------------------------------------
\section{Modelagem Relacional e Django ORM}
%-------------------------------------------------
\label{sec:modelagem-relacional}

A modelagem relacional é a base sobre a qual os sistemas de informação organizam e persistem seus dados.
No contexto do Django, a definição dos modelos determina diretamente a estrutura do banco de dados e,
consequentemente, as possibilidades de consulta e relacionamento entre entidades \cite{elmasri2011}.

Neste projeto, um aspecto relevante da modelagem é a decisão de \textbf{não} vincular o layout de
orçamento diretamente ao modelo \texttt{Orcamento} por meio de uma chave estrangeira. Em vez disso, o
sistema resolve o layout ativo em tempo de geração, consultando o modelo \texttt{LayoutOrcamento} pelo
campo \texttt{is\_active}. Essa decisão arquitetural segue o princípio de baixo acoplamento
(\textit{loose coupling}): o modelo de negócio central (\texttt{Orcamento}) permanece desacoplado das
decisões de apresentação (layout), facilitando a evolução independente de ambos \cite{martin2002}.

O campo \texttt{AutoSlugField}, fornecido pela biblioteca \texttt{django-autoslug}, gera
automaticamente um identificador textual único a partir de outro campo do modelo — neste caso, a
descrição do layout — sem requerer intervenção manual, o que simplifica a identificação de registros em
rotas e na interface administrativa.

%-------------------------------------------------
\section{Arquitetura de Software e Boas Práticas}
%-------------------------------------------------
\label{sec:arquitetura}

O princípio da separação de responsabilidades (\textit{Separation of Concerns}) orienta a organização do
código em camadas com papéis bem definidos \cite{dijkstra1982}. No contexto deste projeto, ele se
manifesta na divisão entre: (a) a camada de modelos, responsável pela persistência e pela lógica de
dados; (b) a camada de \textit{views}, responsável pela orquestração, ou seja, a resolução do layout ativo, montagem
do contexto e acionamento do compilador; e (c) a camada de templates, responsável exclusivamente pela
estrutura visual do documento LaTeX.

Essa separação garante que alterações no layout não exijam modificações na lógica da \textit{view}, e
que mudanças nas regras de negócio não impactem a estrutura do template, reduzindo o custo de
manutenção e aumentando a testabilidade do sistema \cite{martin2002}.